% ============================================================
%  articles/a4-divisors.tex
%  A Classification of Principal Divisors on CP^1
%  — Pablo Santos Guerrero
% ============================================================

\smjarticle{art:divisors}{%
  title       = {A Classification of Principal Divisors on $\mathbb{CP}^1$},
  % cover-title = {},
  % toc-title   = {},
  author      = {Pablo Santos Guerrero},
  year        = {MMDA'27},
  affil       = {UgS V},
  category    = {article},
  abstract    = {Every nonzero polynomial over $\mathbb{C}$ has exactly as many roots
    as its degree. This paper develops the analogue of that statement for the much
    larger class of meromorphic functions on the complex projective line
    $\mathbb{CP}^1 = \mathbb{C} \cup \{\infty\}$, where a function may not only vanish
    but also blow up. To each such function we attach its divisor: the formal record
    of where it vanishes and where it has poles, together with the integer
    multiplicity at each such point. Divisors form an abelian group
    $\operatorname{Div}(\mathbb{CP}^1)$; those arising from functions form a
    subgroup $\operatorname{Prin}(\mathbb{CP}^1)$; and the total multiplicity gives a
    homomorphism $\deg : \operatorname{Div}(\mathbb{CP}^1) \to \mathbb{Z}$. Our main
    result is that on $\mathbb{CP}^1$ these two notions agree exactly:
    $\operatorname{Prin}(\mathbb{CP}^1) = \ker(\deg)$. From this we compute the
    divisor class group $\operatorname{Cl}(\mathbb{CP}^1) \cong \mathbb{Z}$ and,
    by an explicit construction with transition functions, the Picard group
    $\operatorname{Pic}(\mathbb{CP}^1) \cong \mathbb{Z}$.},
}

% Theorem numbering
\counterwithin{thrmctr}{section}
\import{articles/papers/divisors/}{main.tex}
%
\counterwithout{thrmctr}{section}